\section{GM-4. Clase abstracta e interfaz utilizadas}

\subsection*{Aplicación}
\texttt{Gota} es abstracta. Sus hijas \texttt{GotaBuena} y \texttt{ObjetoEspecial} son creadas por \texttt{Lluvia.crearGota()} y usadas como \texttt{Gota} en \texttt{Lluvia.mover()}. La primera recoge agua; la segunda aplica un \texttt{TipoObjeto}. \texttt{GotaMala} sigue en el código como clase de la versión anterior, pero no se utiliza como evidencia principal porque la fábrica actual no la genera.

\texttt{ControlMovimiento} tiene dos implementaciones: \texttt{ControlTeclado} y \texttt{ControlTactil}. \texttt{GameLluvia.iniciar()} elige una y \texttt{Partida} la entrega al personaje. \texttt{Personaje.actualizar()} la invoca en ejecución. El modo ratón permite ejercitar la estrategia táctil en escritorio.

\subsection*{Evidencia}
\begin{itemize}
\item \href{https://github.com/maverick1006/Proyecto-Game-LLuvia-/blob/entrega-final/core/src/main/java/puppy/code/Gota.java#L5}{Gota.java, línea 5} (\texttt{public abstract class})
\item \href{https://github.com/maverick1006/Proyecto-Game-LLuvia-/blob/entrega-final/core/src/main/java/puppy/code/GotaBuena.java#L3}{GotaBuena.java, línea 3} (\texttt{extends Gota})
\item \href{https://github.com/maverick1006/Proyecto-Game-LLuvia-/blob/entrega-final/core/src/main/java/puppy/code/ObjetoEspecial.java#L3}{ObjetoEspecial.java, línea 3} (\texttt{extends Gota})
\item \href{https://github.com/maverick1006/Proyecto-Game-LLuvia-/blob/entrega-final/core/src/main/java/puppy/code/Lluvia.java#L40}{Lluvia.java, línea 40} (\texttt{private Gota crearGota(})
\item \href{https://github.com/maverick1006/Proyecto-Game-LLuvia-/blob/entrega-final/core/src/main/java/puppy/code/ControlMovimiento.java#L3}{ControlMovimiento.java, línea 3} (\texttt{public interface})
\item \href{https://github.com/maverick1006/Proyecto-Game-LLuvia-/blob/entrega-final/core/src/main/java/puppy/code/ControlTeclado.java#L6}{ControlTeclado.java, línea 6} (\texttt{implements ControlMovimiento})
\item \href{https://github.com/maverick1006/Proyecto-Game-LLuvia-/blob/entrega-final/core/src/main/java/puppy/code/ControlTactil.java#L7}{ControlTactil.java, línea 7} (\texttt{implements ControlMovimiento})
\item \href{https://github.com/maverick1006/Proyecto-Game-LLuvia-/blob/entrega-final/core/src/main/java/puppy/code/Personaje.java#L41}{Personaje.java, línea 41} (\texttt{control.actualizar(this, delta)})
\item \href{https://github.com/maverick1006/Proyecto-Game-LLuvia-/blob/entrega-final/core/src/main/java/puppy/code/GameLluvia.java#L50}{GameLluvia.java, línea 50} (\texttt{private void iniciar()})
\end{itemize}

Diseño: diagramas de GM-3, GM-7 y GM-8. No basta declarar las clases: las llamadas anteriores demuestran su uso.

